The central question addressed by this whitepaper is whether a Prajna-driven Grover search combined with a three-tier intelligence stack can constitute a practical architecture for attacking the seven Clay Millennium Prize Problems. Our answer is twofold. In the weak sense, yes: the architecture can accelerate discovery and validation for bounded-oracle fragments of frontier mathematics. In the strong sense, no: no single architecture, classical or quantum, can be a universal solver for open-ended mathematical invention, because the oracle-design problem for frontier conjectures is often as hard as the conjecture itself, because Grover delivers only a quadratic speedup, and because undecidability places a floor under what any finite proof-search system can guarantee.

The contribution of this whitepaper is therefore structural rather than triumphalist. We contribute (i) a first-principles formalisation of LRAM with explicit operators and losses, (ii) a convergence analysis under a bounded-oracle assumption, (iii) an inversion-style enumeration of failure modes, (iv) a higher-dimensional lifting that uses symmetry and categorical structure to compress the candidate space before quantum amplification, (v) a teacher-student distillation recipe that uses large language models, large action models, and small specialist models as complementary teachers, (vi) a Reflexion-enhanced outer loop, and (vii) an end-to-end open-source proof-of-concept and proof-of-value workshop that operationalises every equation in the whitepaper in a reproducible form.

The practical payoff is a research engine for machine-assisted theorem discovery, not a universal oracle for mathematics. The ambitious framing toward the Millennium Prize Problems is preserved and motivates the staged roadmap presented in the final sections.
